% Sólo las referencias citadas en el texto, en estilo IEEE y en orden de aparición.
\begin{thebibliography}{10}

\bibitem{inegi_eda}
Instituto Nacional de Estadística y Geografía, ``Conjunto de datos vectorial edafológico, escala 1:250\,000, Series I
y II, hoja E14-2 Ciudad de México.'' [En línea]. Disponible: \url{https://www.inegi.org.mx/temas/edafologia/}

\bibitem{siap}
Servicio de Información Agroalimentaria y Pesquera, ``Cierre de la producción agrícola por municipio, cebada grano,
temporal, primavera-verano, 2023--2025.'' [En línea]. Disponible: \url{https://nube.agricultura.gob.mx/cierre_agricola/}.
Consultado el 9 de octubre de 2026.

\bibitem{roberts2017}
D.~R. Roberts \emph{et al.}, ``Cross-validation strategies for data with temporal, spatial, hierarchical, or
phylogenetic structure,'' \emph{Ecography}, vol.~40, no.~8, pp.~913--929, 2017.

\bibitem{inifap_hgo}
Instituto Nacional de Investigaciones Forestales, Agrícolas y Pecuarias, ``Agenda técnica agrícola de Hidalgo.''
[En línea]. Disponible:
\url{https://www.gob.mx/cms/uploads/attachment/file/1066180/Agenda_T_cnica_Agr_cola_-_Hidalgo-compressed.pdf}

\bibitem{ceniza_ene25}
``Alertan sobre caída de ceniza en varios municipios de Tlaxcala,'' \emph{Municipios Puebla}, 25 de enero de 2025.
[En línea]. Disponible:
\url{https://archivo.municipiospuebla.mx/nota/2025-01-25/tlaxcala/alertan-sobre-ca%c3%adda-de-ceniza-en-varios-municipios-de-tlaxcala}

\bibitem{ceniza_2025}
Azteca Puebla, reportes diarios de dispersión de ceniza del Popocatépetl, mayo a octubre de 2025. [En línea].
Disponible: \url{https://www.aztecapuebla.com/noticias/}

\bibitem{ligot2023}
N.~Ligot, P.~Bogaert, S.~Biass, G.~Lobet y P.~Delmelle, ``Grain size modulates volcanic ash retention on crop foliage
and potential yield loss,'' \emph{Natural Hazards and Earth System Sciences}, vol.~23, pp.~1355--1369, 2023.

\end{thebibliography}
